% SCRATCH DRAFT — Vol 2 courtroom coda, Ch3 "Orbits and Rotation" (task vol2-coda, author: Podo).
% THIS IS NOT THE BOOK. Do not \input. Floor commit not yet banked; 03.tex untouched.
% CAREFUL CHAPTER — on the ±1 fence.
%
% CONTINUITY: state carries from Ch2 (accepted). Ch2 closed on the wheel about to bend into a loop
% that keeps a residue; the charge made to balance; trace spanning a boundary; momentum the conjugate;
% two coming instruments to agree on a moving number. Ch3 OPENS the loop — the wheel's own truth
% (it goes around) becomes the subject; the count is carried around a closed path and read on return.
%
% PROPOSED SEAM in books/experimentation/latex/chapters/03.tex:
%   KEEP locked body through line 210 ("...how much of what it reports it has seen and how much it has inferred.")
%   KEEP the \rule separator block (lines 212-216).
%   REPLACE the OLD beat now at lines 218-224
%     ("Some accounts do not run in a line from start to finish; they close. ... the account that comes home.")
%   WITH the prose below. The new arc GROWS the old beat's best line
%   ("some accounts do not run in a line, they close --- the account that comes home") into the continuing story.
%
% ±1 FENCE: the residue's VALUE stays entirely out — no "+1", no "minus one", no "exactly one", no
%   committing the residue to a number or to "two" senses. Seed only that the residue has a SENSE and
%   that WHICH way it falls is fixed by the frame WE hold (reverse the frame, the sense reverses), while
%   THAT it exists is the world's. Spinor OUT. Smooth-shadow = this chapter's ceiling, NOT the Ch29 floor.
%
% >>>>>>>>>>>>>>>>>>>>>> APPEND BEGINS (after the \rule) >>>>>>>>>>>>>>>>>>>>>>

\noindent Some accounts do not run in a line; they close. Where the last chapter left the wheel about to bend
back on itself, the record of such a motion no longer runs from one place to another but returns to the very
configuration it began in, coming home to its own first entry period after period. That the path closes, and
keeps closing, is the world's; the mark we set at the start, the point we name the beginning of the round, is
ours. A closed account is a different kind of exhibit from a stretch of road between two places --- a circuit,
the account that comes home --- and the charge the chapter before made to balance now rides one. An orbit is
the plainest of these: a closed account balanced not by a force pinning anything in place but by a
reconciliation that never finishes --- the motion's to keep coming back, the instrument's to keep checking that
it has.

For this was the wheel's own truth, held back until now: it goes around. Every count the sensor laid down was
one pass of a loop, and read straight, in a line, the passes looked like travel forward; but the wheel was
coming back to its start the whole time. So the instrument does the new thing this chapter gives it, the one
every reading after will lean on: it carries a tally around the closed path and reads it not at the end of a
road but where it comes home.

Carry a count around the loop and one of two things happens. Sometimes it returns exactly as it left, nothing
gained on the round --- the tally conserved, the books closing around a cycle just as they closed across a
boundary before, and that unchanged homecoming is a conservation. It comes home unchanged for a reason, not by
luck: a sameness in the motion --- that it reads alike whichever way it is turned --- is what the world keeps,
an invariant bred by that sameness and carried around the loop intact. Sometimes it returns \emph{marked} ---
carrying a residue, a leftover the going-around put there that no single leg of the trip would ever show. That
returning-marked is holonomy, and it is the instrument's decisive new power: to read off the homecoming a thing
that belongs to the whole loop and to none of its parts. (The smooth circle the instrument draws through the
loop's counts is, as ever, a shadow --- the shape the discrete round approaches as the passes are imagined
refined, ours to sketch and never a count the instrument holds. It is this chapter's honest ceiling, named as a
shadow and not pressed past what it is.)

The residue asks to be handled with care, because a line runs through it that the book does not step over yet.
The residue has a \emph{sense} --- a which-way it leaves the tally facing when the loop comes home. Carry not a
number this time but a direction around the closed path, holding it as straight as the path will allow and
never turning it by hand, and it comes home rotated all the same: the rule of ``as straight as allowed'' is
ours to apply, the bend the loop forces on it the world's, and the turn is the residue made visible. But which
way that sense falls does not belong to the world; it is fixed only once we lay down a frame to read it against, an
orientation we call ``unturned'' and hold still by hand. Read the homecoming against that frame and the sense
is fixed with it; read it against the reverse and the sense reverses, though the loop wound exactly as it wound
and the world did nothing. \emph{That} the loop leaves a residue at all is the world's --- the closing put it
there, and no framing of ours talks it away. \emph{Which} way the residue points waits on a hand to give the
loop a direction. The book marks the residue as real here and no more; the sense it settles to is reserved for
where the instrument can make it stand.

Here, too, the book's oldest identity comes true to the letter. From the first page the charge was the count;
now the count is exactly what the loop carries around and brings back --- the charge \emph{is} the loop's own
tally, gone all the way around and returned. Charge equals what the loop counts: the phrase was always this.
And the trace, which in the last chapter spanned a boundary, is carried one link further --- it is now what the
loop retains of the turning across a full circuit, the accumulated rotation brought the whole way around and
read where it comes home, residue and all. The charge borne on the record is that loop count; the payment, when
it is finally named, will be nothing but the loop count read to its end. The chapter need not spend it; it is
enough that the going-around is now, plainly, the thing being counted and the thing that will be owed.

A tally carried around a closed loop is still not the record on one instrument's word. The car's own reading,
gone around and come home, will have to be met by the two instruments named as coming and made to agree on the
same count --- residue and sense and all, three independent hands brought to one closed number, not merely one.
Their physics still waits for the parts that carry it; the demand only tightens, that a charge which goes around
must come home to the same value in every hand that reads it. And the record that must close is not closed. Part
I ends next with the stretches the instrument did not watch end to end --- where it reconstructs the motion it
did not directly record and says, honestly, how much of what it reports it saw and how much it inferred. The
loop is entered; the account has come home once, but the case has not. The court does not sit until the record
is complete.

% <<<<<<<<<<<<<<<<<<<<<< APPEND ENDS <<<<<<<<<<<<<<<<<<<<<<
